\chapter{Estruturas de Texto e Tabelas}

Este capítulo funciona como uma página de referência para o desenho interno de livros técnicos da Versatus HPC. O conteúdo é propositalmente genérico: ele deve ser substituído pelo texto real do projeto, mas preserva exemplos de hierarquia, ritmo editorial, tabelas e parágrafos contínuos. A intenção é facilitar a validação do template antes do início da escrita final.

\section{Três parágrafos contínuos}

A documentação técnica deve privilegiar uma sequência de leitura estável: contexto, decisão, justificativa e consequência. Em materiais didáticos, essa ordem reduz ambiguidade e evita que o leitor precise inferir a arquitetura mental do autor. O texto corrido deve ser usado quando a explicação depende de continuidade conceitual, especialmente em capítulos introdutórios, fundamentos teóricos e discussões de projeto.

Em documentos de engenharia, parágrafos longos demais tendem a prejudicar a manutenção editorial. A recomendação é preservar blocos com uma ideia central por parágrafo, usando listas apenas quando houver enumeração real. Essa disciplina ajuda a converter o mesmo conteúdo em livro, manual, relatório técnico ou white paper sem reescrita estrutural excessiva.

Para programas de estudo, a redação deve separar claramente conceitos estáveis, hipóteses de trabalho, exemplos, restrições e decisões normativas. Essa separação permite que o material seja usado tanto para leitura sequencial quanto para consulta posterior. O template foi desenhado para aceitar ambos os modos sem alterar a identidade visual.

\section{Tabela pequena}

A tabela pequena deve ser usada para parâmetros, convenções, nomenclaturas ou comparações curtas. Ela não deve dominar a página.

\begin{table}[htbp]
\centering
\caption{Exemplo de tabela pequena}
\label{tab:pequena}
\begin{tabular}{lll}
\toprule
Item & Valor & Observação \\
\midrule
CPU & 2 sockets & Nós de computação \\
GPU & 8 unidades & Treinamento de modelos \\
Rede & 400 Gb/s & Interconexão de baixa latência \\
\bottomrule
\end{tabular}
\end{table}

\section{Tabela média}

A tabela média é indicada quando há mais texto por célula. Neste caso, \texttt{tabularx} permite que as colunas se ajustem à largura útil da página.

\begin{table}[htbp]
\centering
\caption{Exemplo de tabela média com largura controlada}
\label{tab:media}
\begin{tabularx}{\textwidth}{>{\bfseries}p{0.22\textwidth}X X}
\toprule
Categoria & Uso recomendado & Critério de aceite \\
\midrule
Provisionamento & Descrever instalação, imagem base, dependências e sequência de boot. & O leitor consegue reproduzir o ambiente sem decisões implícitas. \\
Operação & Explicar rotinas administrativas, telemetria e tratamento de falhas. & O procedimento apresenta pré-condições, passos e rollback. \\
Validação & Registrar testes, métricas, hipóteses e limites conhecidos. & Os resultados são rastreáveis e repetíveis. \\
\bottomrule
\end{tabularx}
\end{table}

\section{Tabela grande}

A tabela grande deve ser usada com moderação. Em livros e manuais, ela é útil para matrizes de requisitos, listas de componentes, parâmetros de configuração e inventários técnicos. Quando a tabela atravessar páginas, use \texttt{longtable}.

{\scriptsize
\setlength{\tabcolsep}{3pt}
\begin{longtable}{p{0.08\textwidth}p{0.15\textwidth}p{0.14\textwidth}p{0.12\textwidth}p{0.16\textwidth}p{0.14\textwidth}p{0.11\textwidth}}
\caption{Exemplo de tabela grande com múltiplas linhas}\label{tab:grande}\\
\toprule
ID & Componente & Função & Prioridade & Métrica & Evidência & Estado \\
\midrule
\endfirsthead
\toprule
ID & Componente & Função & Prioridade & Métrica & Evidência & Estado \\
\midrule
\endhead
\midrule
\multicolumn{7}{r}{Continua na próxima página}\\
\endfoot
\bottomrule
\endlastfoot
R01 & Kernel & Base de execução do sistema & Alta & Latência e estabilidade & Testes de carga & Aberto \\
R02 & Rede & Comunicação entre nós & Alta & Vazão e perda de pacotes & Benchmarks sintéticos & Aberto \\
R03 & Storage & Persistência paralela & Alta & IOPS e throughput & IOR, mdtest & Aberto \\
R04 & Scheduler & Alocação de recursos & Alta & Tempo de fila & Simulações de workload & Aberto \\
R05 & Observabilidade & Eventos e métricas & Média & Custo de coleta & Telemetria contínua & Aberto \\
R06 & Segurança & Controle de acesso & Alta & Superfície exposta & Revisão de configuração & Aberto \\
R07 & Documentação & Reprodutibilidade & Média & Cobertura editorial & Checklist de publicação & Aberto \\
R08 & CI/CD & Geração de artefatos & Média & Tempo de build & Pipeline automatizado & Aberto \\
R09 & Instalação & Primeira implantação & Alta & Taxa de sucesso & Testes em ambiente limpo & Aberto \\
R10 & Atualização & Evolução controlada & Média & Rollback funcional & Ensaios de versão & Aberto \\
R11 & Diagnóstico & Suporte operacional & Média & Tempo de triagem & Casos simulados & Aberto \\
R12 & Treinamento & Transferência de conhecimento & Baixa & Conclusão dos módulos & Avaliação interna & Aberto \\
\end{longtable}
}
